\scanpage{214}
% slanted
\margin{\uline{Scholie}}

En résumé, nous \deleted{\unsure{sommes ainsi|sommes ici}} amenés à nous intéresser à la structure suivante

\[
\left\{\begin{aligned}
&\text{a) Triple }\Gamma=(S,S^{\unsure{*|\vee}},R)\ \text{de deux ens }S,S^\vee\ \text{et d'une}\\
&\qquad \text{relation }R\subset S\times S^\vee\\
&\text{b) Donnée d'un ss-groupe \uline{commutatif}}\\
&\hspace{20mm}Z\subset\operatorname{Aut}(\Gamma),\\
&\qquad\text{dont le centralisateur }\inserted{\text{dans }\operatorname{Aut}(\Gamma)}\ \text{est noté }\unsure{\mathcal G|G}.
\end{aligned}\right. \tag*{(69)}
\]

% slanted
\margin{dans le cas de l'exemple linéaire avec $M^*,M^{\vee *}\ldots$, $\mathcal G$ contient en tous cas le groupe $\operatorname{Gl}_\Lambda\bigl(\mathcal O(\unsure{*,M|M})\bigr)'$}

On appelle \uline{$Z$-jauges} d'une telle structure les \deleted{applications} bijections

\[
\left\{\begin{aligned}
&\varphi:S\xrightarrow{\sim} S^{\unsure{*|\vee}}\qquad\text{satisfaisant}\\
&1)\quad\deleted{\unsure{\varphi(s,t)}}(\mathrm{id}\times\varphi)^{-1}(R)\ \text{est une relation \uline{symétrique}}\\
&\qquad\text{sur }S\\
&\qquad\text{i.e. }\forall s,t\in S,\quad(s,\varphi(t))\in R\iff(t,\varphi(s))\in R\\
&2)\quad\forall(\lambda,\lambda^\vee)\in Z,\ \text{on a}\\
&\qquad\lambda^\vee\varphi\lambda=\varphi\qquad\text{i.e }\lambda^\vee\varphi=\varphi\lambda^{-1}
\end{aligned}\right. \tag*{(70)}
\]

Si $R$ est biséparante, de sorte que \deleted{$S^\vee$ soit} la donnée de $R$ s'identifie à celle de $(S,S^\vee\subset\mathscr P(S))$, alors les $\varphi$ satisfaisant (1) s'identifient aux structures de graphe \deleted{\unsure{sép}} \inserted{strict} séparant sur $S$ (sans exclure des boucles) $A_0\subset\mathscr P_2(S)$, qui admettent $S^\vee$ comme ens de leurs étoiles, et pour lesquels les $\lambda\in Z_S$ (provenant des $\lambda\in Z$) sont des \uline{homothéties}. \inserted{En tous cas,} Les $\varphi$ en question, \inserted{(s'il y en a !)}, forment un \uline{torseur} (non sous $Z$, mais) sous $\mathcal G=\operatorname{Cent}_{\operatorname{Aut}(\Gamma)}(Z)$.

Les $Z$-multigraphes envisagés précédemment sont donc les sous \inserted{$Z$-}torseurs de celui-ci, \deleted{i.e des} \inserted{donc s'identifient} aux points de $\mathcal G/Z$.

% slanted
\margin{canular, cf annotation page précédente. Mais $Z$ opère librement sur $\Pi_Z$}

% archive sheet number: 213
